\documentclass[11pt]{article}

\usepackage[margin=1in]{geometry}
\usepackage[T1]{fontenc}
\usepackage[utf8]{inputenc}
\usepackage{lmodern}
\usepackage{xcolor}
\usepackage{titlesec}
\usepackage{fancyhdr}
\usepackage{enumitem}
\usepackage{booktabs}
\usepackage{tabularx}
\usepackage{graphicx}
\usepackage{caption}
\usepackage{listings}
\usepackage[most]{tcolorbox}
\usepackage{abstract}
\usepackage[hidelinks]{hyperref}

\definecolor{accent}{HTML}{1F3A5F}
\definecolor{accentgray}{HTML}{F3F5F8}
\definecolor{linegray}{HTML}{D4D9E1}
\definecolor{textgray}{HTML}{333333}

\hypersetup{
  colorlinks=true,
  linkcolor=accent,
  urlcolor=accent,
  citecolor=accent,
  pdftitle={ContractLens Milestone 2 Architecture Report},
  pdfauthor={ContractLens Team},
  pdfsubject={Milestone 2 Architecture}
}

\titleformat{\section}
  {\large\bfseries\color{accent}}
  {\thesection.}
  {0.5em}
  {}
\titlespacing*{\section}{0pt}{14pt}{8pt}

\titleformat{\subsection}
  {\normalsize\bfseries\color{accent}}
  {\thesubsection}
  {0.5em}
  {}
\titlespacing*{\subsection}{0pt}{10pt}{6pt}

\setlength{\parskip}{6pt}
\setlength{\parindent}{0pt}
\setlist[itemize]{leftmargin=1.5em, itemsep=2pt, topsep=4pt}
\setlist[enumerate]{leftmargin=1.7em, itemsep=2pt, topsep=4pt}

\pagestyle{fancy}
\fancyhf{}
\fancyhead[L]{\textcolor{accent}{\footnotesize ContractLens MS2 Report}}
\fancyhead[R]{\textcolor{accent}{\footnotesize \nouppercase{\leftmark}}}
\fancyfoot[C]{\footnotesize \thepage}
\renewcommand{\headrulewidth}{0.3pt}
\renewcommand{\footrulewidth}{0pt}

\fancypagestyle{plain}{
  \fancyhf{}
  \fancyhead[L]{\textcolor{accent}{\footnotesize ContractLens MS2 Report}}
  \fancyhead[R]{\textcolor{accent}{\footnotesize \nouppercase{\leftmark}}}
  \fancyfoot[C]{\footnotesize \thepage}
  \renewcommand{\headrulewidth}{0.3pt}
}

\lstdefinelanguage{json}{
  basicstyle=\ttfamily\small,
  numbers=left,
  numberstyle=\tiny\color{gray},
  stepnumber=1,
  numbersep=8pt,
  showstringspaces=false,
  breaklines=true,
  frame=single,
  rulecolor=\color{linegray},
  backgroundcolor=\color{accentgray},
  keywordstyle=\color{accent},
  stringstyle=\color{teal!60!black},
  commentstyle=\color{gray},
  morestring=[b]\"
}

\newtcolorbox{abstractbox}{
  colback=accentgray,
  colframe=linegray,
  boxrule=0.5pt,
  arc=2pt,
  left=8pt,
  right=8pt,
  top=8pt,
  bottom=8pt
}

\begin{document}

\begin{titlepage}
  \thispagestyle{empty}
  \centering
  \vspace*{2.5cm}
  {\Huge\bfseries\color{accent} ContractLens Milestone 2 Architecture Report\par}
  \vspace{0.8cm}
  {\Large Milestone 2 --- Agentic NDA Review with Retrieval-Augmented Traceability\par}
  \vspace{1.8cm}
  {\large Milestone 2\par}
  \vspace{0.5cm}
  {\large ContractLens / ContractNLI Project\par}
  \vspace{0.5cm}
  {\large \today\par}
  \vspace{1.0cm}
  {\large Team Attribution: ContractLens Milestone Team\par}
  \vspace{0.35cm}
  \begin{minipage}{0.72\textwidth}
    \centering
    \small
    \begin{tabularx}{\textwidth}{>{\raggedright\arraybackslash}X >{\raggedright\arraybackslash}p{2.2cm}}
      \toprule
      \textbf{Name} & \textbf{ID} \\
      \midrule
      Daniel Ashraf Ekdawi & 55-1598 \\
      Mostafa Ashraf Ezzat & 55-5414 \\
      Ahmed Hatim Backar & 55-24857 \\
      David George Atef & 55-26933 \\
      Youssef Mohamed Malek & 55-24816 \\
      \bottomrule
    \end{tabularx}
  \end{minipage}
  \vfill
  \begin{minipage}{0.86\textwidth}
    \small
    \textbf{Executive Summary.} This report presents the Milestone 2 architecture of ContractLens, extending a batch legal-NLI baseline into an agentic, conversational NDA review system. The architecture combines orchestration, retrieval support, and schema-governed traceability while maintaining strict legal grounding constraints. A central design principle is preserved throughout: retrieval examples inform reasoning, but final legal evidence must come exclusively from the analyzed contract.
  \end{minipage}
  \vspace*{1.3cm}
\end{titlepage}

\tableofcontents
\newpage

\begin{abstractbox}
\textbf{Abstract.} This report documents the Milestone 2 architecture of ContractLens, an NDA review system built on ContractNLI that evolves from batch legal-NLI evaluation into an agentic and conversational workflow. The design combines \texttt{Qwen3-4B} orchestration, a fine-tuned \texttt{Qwen3-1.7B} NLI core, retrieval support through Vector RAG and GraphRAG, hypothesis-level decomposition, and schema-driven auditability through updated playbook metadata and \texttt{RunTrace 2.0}. The report emphasizes a core safety principle: retrieved training examples are external reasoning memory only, while final legal evidence must be grounded exclusively in spans from the analyzed contract.
\end{abstractbox}

\section{Introduction / Project Context}
ContractLens is an NDA review system built around the ContractNLI task setting, where fixed legal hypotheses are evaluated against contract text with evidence grounding. In Milestone 1, the project delivered a batch legal-NLI baseline that focused on per-contract hypothesis classification and trace generation in a mostly linear pipeline.

Milestone 2 extends that baseline into an agentic and conversational architecture. The system is designed to support user-driven contract questioning while preserving legal grounding constraints and auditable decision traces.

The main Milestone 2 additions are: \texttt{Qwen3-4B} orchestration, vector and graph retrieval branches, conversation memory, playbook metadata for per-hypothesis agent attribution, and \texttt{RunTrace 2.0} auditability features that separate \texttt{external\_memory} from \texttt{contract\_evidence}.

\section{System Overview and Design Rationale}
Milestone 1 was a batch legal-NLI baseline: a single fine-tuned model processed one contract and produced hypothesis labels/evidence in a single pipeline. Milestone 2 evolves this into an agentic NDA review architecture with explicit decomposition, retrieval-assisted reasoning, and traceable component boundaries.

The target execution pattern is \textbf{broadcast $\rightarrow$ fan-out $\rightarrow$ aggregate}: the NLI core performs a single broad pass, task-specific hypothesis work is decomposed into units, and a final aggregation stage materializes a unified trace artifact.

This decomposition improves traceability (per-hypothesis provenance and metadata), modularity (independent RAG branches and workers), and maintainability (clear separation of orchestration, retrieval, inference, and aggregation concerns). Table~\ref{tab:component-summary} summarizes the primary components and their current status.

\begin{table}[htbp]
\centering
\caption{Component Summary}
\label{tab:component-summary}
\small
\begin{tabularx}{\textwidth}{>{\raggedright\arraybackslash}p{2.2cm} >{\raggedright\arraybackslash}p{2.3cm} >{\raggedright\arraybackslash}X >{\raggedright\arraybackslash}p{2.6cm}}
\toprule
\textbf{Component} & \textbf{Model / Method} & \textbf{Responsibility} & \textbf{Implementation Status} \\
\midrule
Orchestrator & \texttt{Qwen3-4B} & Session-level routing, tool-calling, path selection between conversation and analytical flow. & Architectural / MS3-target \\
NLI Core & Fine-tuned \texttt{Qwen3-1.7B} & Single-pass evaluation of all 17 hypotheses with structured outputs and evidence spans. & Present \\
Vector RAG & Embedding + FAISS & Semantic retrieval from \texttt{data/train.json} for reasoning support. & Architectural / pending \\
GraphRAG & Concept graph retrieval & Concept- and hypothesis-targeted retrieval over graph schema. & Architectural / pending \\
Conversation Agent & \texttt{Qwen3-4B} + history & User-facing Q\&A, context injection, conversation continuity. & Present (\texttt{src/conversation\_agent.py}, \texttt{agent.py}) \\
Playbook & Schema-governed ruleset & Policy mapping from hypothesis outcomes to severity/action/rationale. & Present (schema updated) \\
\texttt{RunTrace 2.0} & Schema-governed trace model & Auditable run output with retrieval/evidence separation and agent metadata. & Present (schema updated) \\
\bottomrule
\end{tabularx}
\end{table}

\IfFileExists{architecture/architecture.pdf}{
\begin{figure}[htbp]
  \centering
  \includegraphics[width=0.9\textwidth]{architecture/architecture.pdf}
  \caption{ContractLens Milestone 2 architecture overview}
  \label{fig:arch-overview}
\end{figure}
}{
\noindent\textit{Architecture figure omitted: \texttt{architecture/architecture.pdf} not found.}
}

As shown in Figure~\ref{fig:arch-overview}, orchestration and retrieval are separated from evidence-grounding logic, reducing the risk of retrieval leakage into final legal assertions.

This architectural framing leads directly to the orchestrator's role as routing and control logic, which coordinates model inference, retrieval, and output structuring.

\section{Orchestrator}
The orchestrator is specified as \texttt{Qwen3-4B} with thinking enabled. It is the tool-calling controller that routes execution across conversation handling and analytical tools.

Architecture specifies JSON-style function calling with call shape:
\texttt{\{"name": "<tool\_name>", "arguments": \{...\}\}}

Planned orchestration routes include:
\begin{itemize}
  \item Conversation path: user prompt + history + contract + retrieval context to conversational answer.
  \item Analytical path: NLI core to dispatcher/task fan-out to hypothesis workers to aggregation.
  \item Retrieval routing: selects exactly one branch per run (\texttt{vector} or \texttt{graph}) and passes mode through downstream traces.
\end{itemize}

The orchestrator is architectural in \texttt{architecture/architecture.yaml} (MS3-target end-to-end control loop).

Because orchestration quality depends on stable structured outputs from the legal classifier, the NLI core is constrained to a dedicated model mode and output discipline.

\section{NLI Core}
The NLI core is the fine-tuned \texttt{Qwen3-1.7B} adapter model:
\begin{quote}
\texttt{Youssef-Malek/}\\
\texttt{contractnli-vast-ai-qwen3-1.7b}
\end{quote}
It is designed to run once per contract over all 17 hypotheses.

Key policy:
\begin{itemize}
  \item \textbf{Thinking OFF} for the fine-tuned NLI core (\texttt{enable\_thinking=False}).
  \item Rationale: the adapter was trained for direct structured classification output; enabling thinking introduces format drift risk.
\end{itemize}

Expected NLI core output is hypothesis-level classification with contract-grounded evidence span indices (ENTAILED, CONTRADICTED, NOT\_MENTIONED + evidence spans).

After NLI classification, retrieval is used to support reasoning and calibration, but retrieval output remains separate from final legal evidence.

\section{End-to-End Data Flow}
The intended Milestone 2 end-to-end flow is presented below. The sequence clarifies the boundary between reasoning support and admissible legal grounding:
\begin{enumerate}
  \item The user submits a contract and a question.
  \item The conversation agent loads the analyzed contract.
  \item The orchestrator selects either vector retrieval or graph retrieval.
  \item The selected retrieval module fetches examples from \texttt{data/train.json} only.
  \item Retrieved examples are injected as \texttt{external\_memory}, not final evidence.
  \item The fine-tuned \texttt{Qwen3-1.7B} NLI core evaluates all 17 hypotheses.
  \item Hypothesis tasks are enriched with playbook decisions and metadata.
  \item The aggregator merges outputs into \texttt{RunTrace 2.0}.
  \item The conversation agent returns a grounded answer to the user.
\end{enumerate}

A central rule applies across all steps: final legal evidence must come from spans in the analyzed contract, never from retrieved training examples.

\section{Vector RAG}
Vector RAG is defined to index and retrieve from \texttt{data/train.json} only.

Design details:
\begin{itemize}
  \item Embedding model: \texttt{sentence-transformers/all-MiniLM-L6-v2}.
  \item Similarity index: FAISS \texttt{IndexFlatIP}.
  \item Vectors are L2-normalized so inner product approximates cosine similarity.
  \item Supports both free-text retrieval and hypothesis-targeted retrieval (with optional label-aware filtering/re-ranking).
\end{itemize}

\texttt{data/test.json} is evaluation-only and must not be indexed.

Vector retrieval provides high-coverage semantic matching, while the graph branch complements it with more explicit concept-level interpretability.

\section{GraphRAG}
GraphRAG is defined with three node types:
\begin{itemize}
  \item \texttt{SpanNode}
  \item \texttt{ConceptNode}
  \item \texttt{HypothesisNode}
\end{itemize}

And three edge types:
\begin{itemize}
  \item \texttt{CONTAINS}
  \item \texttt{INVOLVES}
  \item \texttt{CITED\_FOR}
\end{itemize}

The design uses hardcoded legal concept vocabularies and synonyms, intentionally avoiding external NLP/NER dependencies. Retrieval supports concept-based free-text graph traversal and hypothesis-targeted lookup paths.

With retrieval context prepared, the conversation layer can answer user queries while enforcing grounding and provenance constraints.

\section{Conversation Agent}
The conversation agent is user-facing and uses \texttt{Qwen3-4B} with thinking enabled. In the repository, \texttt{src/conversation\_agent.py} and \texttt{agent.py} implement the CLI interaction loop and persist history in \texttt{conversation\_history.json}.

Behavioral constraints:
\begin{itemize}
  \item Retrieval context is injected as \textbf{external memory} only.
  \item Retrieved training examples must never be cited as final legal evidence.
  \item Final legal conclusions and quotes must be grounded in the currently analyzed contract.
\end{itemize}

The implementation enforces this distinction in prompt instructions and structured trace design.

The schema updates in Milestone 2 formalize these runtime requirements into explicit, auditable data structures.

\section{Updated Playbook}
Milestone 2 updates the playbook schema while preserving Milestone 1 legal logic:
\begin{itemize}
  \item Added \texttt{schema\_version: "2.0-ms2"}
  \item Added run-level \texttt{retrieval\_mode} (\texttt{vector} or \texttt{graph})
  \item Added per-check \texttt{agent\_id}
  \item Added per-check \texttt{conversation\_turn} (nullable integer)
\end{itemize}

These additions provide the metadata needed for agentic execution and traceability without changing the underlying legal risk semantics from Milestone 1.

Illustrative MS2 playbook shape:
\begin{lstlisting}[language=json]
{
  "schema_version": "2.0-ms2",
  "retrieval_mode": "vector",
  "checks": {
    "H06": {
      "agent_id": "H06",
      "conversation_turn": null,
      "criticality": "HIGH"
    }
  }
}
\end{lstlisting}

Unchanged elements:
\begin{itemize}
  \item 17 hypothesis rules (H01-H17)
  \item Criticality tiers
  \item Override mappings
  \item Rationale templates
  \item Label/severity/action logic
\end{itemize}

\section{Updated RunTrace Schema}
Milestone 2 extends \texttt{RunTrace 2.0} to capture agentic execution context:
\begin{itemize}
  \item Added \texttt{schema\_version: "2.0-ms2"}
  \item Added \texttt{retrieval\_mode}
  \item Added \texttt{conversation\_history}
  \item Added \texttt{retrieval\_context}
  \item Added per-hypothesis \texttt{agent\_metadata}
\end{itemize}

Illustrative MS2 RunTrace shape:
\begin{lstlisting}[language=json]
{
  "schema_version": "2.0-ms2",
  "retrieval_mode": "graph",
  "retrieval_context": {
    "external_memory": [],
    "contract_evidence": []
  },
  "hypothesis_traces": [
    {
      "hypothesis_id": "H06",
      "agent_metadata": {
        "agent_id": "H06",
        "rag_query": "consultant disclosure",
        "rag_hits": 5,
        "rag_mode": "graph"
      }
    }
  ]
}
\end{lstlisting}

Most importantly, the schema explicitly separates \texttt{external\_memory} and \texttt{contract\_evidence}. Table~\ref{tab:evidence-boundary} summarizes the evidence policy boundary used throughout this report.

\begin{table}[htbp]
\centering
\caption{Evidence Boundary Policy}
\label{tab:evidence-boundary}
\small
\begin{tabularx}{\textwidth}{>{\raggedright\arraybackslash}p{3.2cm} >{\raggedright\arraybackslash}X >{\raggedright\arraybackslash}p{3.8cm}}
\toprule
\textbf{Source} & \textbf{Used For} & \textbf{May Be Cited as Final Legal Evidence?} \\
\midrule
Retrieved \texttt{train.json} examples (\texttt{external\_memory}) & Reasoning support, calibration, comparative context & No \\
Analyzed contract spans (\texttt{contract\_evidence}) & Grounded legal justification and quoted evidence & Yes \\
Conversation history & Context continuity across turns & No \\
NLI labels & Classification output and routing cues & Only with supporting contract evidence \\
\bottomrule
\end{tabularx}
\end{table}

Final decision evidence is constrained to contract evidence shape to preserve legal grounding integrity.

\section{Validation and Safety Constraints}
The Milestone 2 architecture and schemas enforce the following constraints:
\begin{itemize}
  \item Only \texttt{data/train.json} may be indexed for retrieval.
  \item \texttt{data/test.json} is evaluation-only and must not be indexed.
  \item Retrieved examples are external memory only and are not valid final legal evidence.
  \item Final evidence must map to spans in the analyzed contract.
  \item NLI outputs and trace artifacts must remain schema-valid.
  \item Quote integrity and groundedness checks should be preserved.
  \item \texttt{Qwen3-4B} orchestrator/conversation may use thinking.
  \item Fine-tuned \texttt{Qwen3-1.7B} NLI core must run with thinking disabled to avoid structured-output drift.
\end{itemize}

These constraints complement the evidence boundary shown in Table~\ref{tab:evidence-boundary} and act as the operational safety envelope for Milestone 2 behavior.

\section{Limitations and Future Work}
Current design limitations and next steps include:
\begin{itemize}
  \item Vector RAG and GraphRAG are separate retrieval modes, not a hybrid retrieval strategy.
  \item Vector retrieval depends on embedding quality and can surface semantically similar but legally weak examples.
  \item Graph retrieval is interpretable but constrained by manually defined concept and synonym coverage.
  \item Long contracts can create context-window pressure for end-to-end orchestration and conversation.
  \item The conversation layer must continuously avoid citing retrieved training examples as legal evidence.
\end{itemize}

Architectural and pending work items:
\begin{itemize}
  \item Dispatcher implementation (architecturally specified).
  \item Hypothesis worker pool implementation (architecturally specified).
  \item Production aggregator implementation (architecturally specified).
  \item Quote repair and post-processing improvements.
  \item Hybrid retrieval research (vector + graph orchestration).
  \item UI/API integration beyond current CLI-oriented interfaces.
\end{itemize}

\section{Implemented vs Architectural / MS3}
The implementation split below keeps status statements file-truthful while distinguishing present artifacts from planned MS3 orchestration components.

\begin{table}[htbp]
\centering
\caption{Implementation Status Summary}
\label{tab:impl-status}
\small
\begin{tabularx}{\textwidth}{>{\raggedright\arraybackslash}p{4.0cm} >{\raggedright\arraybackslash}p{2.6cm} >{\raggedright\arraybackslash}X}
\toprule
\textbf{Artifact} & \textbf{Current Status} & \textbf{Notes} \\
\midrule
\texttt{architecture/architecture.yaml} & Present & Living architecture specification. \\
\texttt{architecture/architecture.pdf} & Present & Diagram deliverable from architecture spec. \\
\texttt{src/types.py} & Present & Shared Milestone 2 TypedDict contracts. \\
\texttt{src/model\_loader.py} & Present & Device and model-loading logic. \\
\texttt{src/conversation\_agent.py}, \texttt{agent.py} & Present & Conversation flow and CLI entry point available. \\
\texttt{src/rag\_vector.py} & Not present & Architectural/pending (Member 2 deliverable). \\
\texttt{src/rag\_graph.py} & Not present & Architectural/pending (Member 3 deliverable). \\
\texttt{03\_build\_index.py} & Not present & Architectural/pending shared indexing entry point. \\
Dispatcher / worker pool / aggregator modules & Not present & Architectural/MS3-target orchestration components. \\
\bottomrule
\end{tabularx}
\end{table}

As reflected in Table~\ref{tab:impl-status}, Vector RAG and GraphRAG are architecturally specified and expected as Member 2/3 deliverables, while full dispatcher/worker/aggregator orchestration remains MS3-target unless later added.

\section{Conclusion}
Milestone 2 transforms ContractLens from a batch legal-NLI evaluator into an agentic, conversational NDA review system that combines a \texttt{Qwen3-4B} orchestrator, a fine-tuned \texttt{Qwen3-1.7B} NLI core, vector retrieval, graph retrieval, hypothesis-level decomposition, an updated playbook, and \texttt{RunTrace 2.0} auditability.

The central design principle remains strict separation between \texttt{external\_memory} and final contract-grounded legal evidence via \texttt{contract\_evidence}. This principle preserves legal traceability and prevents retrieval support artifacts from being misrepresented as contractual facts.

\end{document}
